{

\#include <stdio.h>}





\bigskip



{

\#include {\textquotedbl}BPatch.h{\textquotedbl}}



{

\#include {\textquotedbl}BPatch_addressSpace.h{\textquotedbl}}



{

\#include {\textquotedbl}BPatch_process.h{\textquotedbl}}



{

\#include {\textquotedbl}BPatch_binaryEdit.h{\textquotedbl}}



{

\#include {\textquotedbl}BPatch_point.h{\textquotedbl}}



{

\#include {\textquotedbl}BPatch_function.h{\textquotedbl}}





\bigskip



{

using namespace std;}



{

using namespace Dyninst;}





\bigskip



{

// Create an instance of class BPatch}



{

BPatch bpatch;}





\bigskip



{

// Different ways to perform instrumentation}



{

typedef enum {}



{

create,}



{

attach,}



{

open}



{

} accessType_t;





\bigskip



{

// Attach, create, or open a file for rewriting}



{

BPatch_addressSpace* startInstrumenting(accessType_t accessType,}



{

const char* name,}



{

int pid,}



{

const char* argv[]) {}



{

BPatch_addressSpace* handle = NULL;}





\bigskip



{

switch(accessType) {}



{

case create:}



{

handle = bpatch.processCreate(name, argv);}



{

if (!handle) { fprintf(stderr, {\textquotedbl}processCreate

failed{\textbackslash}n{\textquotedbl}); }}



{

break;}



{

case attach:}



{

handle = bpatch.processAttach(name, pid);}



{

if (!handle) { fprintf(stderr, {\textquotedbl}processAttach

failed{\textbackslash}n{\textquotedbl}); }}



{

break;}





\bigskip



{

case open:}



{

// Open the binary file; do not open dependencies}



{

handle = bpatch.openBinary(name, false);}



{

if (!handle) { fprintf(stderr, {\textquotedbl}openBinary

failed{\textbackslash}n{\textquotedbl}); }}



{

break;}



{

}}





\bigskip



{

return handle;}



{

}}





\bigskip



{

bool instrumentMemoryAccesses(BPatch_addressSpace* app) {}



{

BPatch_image* appImage = app->getImage();}





\bigskip



{

// We're interested in loads and stores}



{

BPatch_Set<BPatch_opCode> axs;}



{

axs.insert(BPatch_opLoad);}



{

axs.insert(BPatch_opStore);}





\bigskip



{

// Scan the function InterestingProcedure



{

// and create instrumentation points}



{

std::vector<BPatch_function*> functions;}



{

appImage->findFunction({\textquotedbl}InterestingProcedure{\textquotedbl}, functions);}



{

std::vector<BPatch_point*>* points =



{

functions[0]->findPoint(axs);}



{

if (!points) {}



{

fprintf(stderr, {\textquotedbl}No load/store points found{\textbackslash}n{\textquotedbl});}



{

return false;}



{

}}





\bigskip



{

// Create the printf function call snippet}



{

std::vector<BPatch_snippet*> printfArgs;}



{

BPatch_snippet* fmt = new BPatch_constExpr({\textquotedbl}Access at: 0x\%lx{\textbackslash}n{\textquotedbl});}



{

printfArgs.push_back(fmt);}



{

BPatch_snippet* eae = new BPatch_effectiveAddressExpr();}



{

printfArgs.push_back(eae);}





\bigskip



{

// Find the printf function}



{

std::vector<BPatch_function*> printfFuncs;}



{

appImage->findFunction({\textquotedbl}printf{\textquotedbl}, printfFuncs);}



{

if (printfFuncs.size() == 0) {}



{

fprintf(stderr, {\textquotedbl}Could not find printf{\textbackslash}n{\textquotedbl});}



{

return false;}



{

}}





\bigskip



{

// Construct a function call snippet



{

BPatch_funcCallExpr printfCall(*(printfFuncs[0]), printfArgs);}





\bigskip



{

// Insert the snippet at the instrumentation points}



{

if (!app->insertSnippet(printfCall, *points)) {}



{

fprintf(stderr, {\textquotedbl}insertSnippet failed{\textbackslash}n{\textquotedbl});}



{

return false;}



{

}}



{

return true;}



{

}}





\bigskip



{

void finishInstrumenting(BPatch_addressSpace* app, const char* newName) {}



{

BPatch_process* appProc = dynamic_cast<BPatch_process*>(app);}



{

BPatch_binaryEdit* appBin = dynamic_cast<BPatch_binaryEdit*>(app);}





\bigskip



{

if (appProc) {}



{

if (!appProc->continueExecution()) {}



{

fprintf(stderr, {\textquotedbl}continueExecution failed{\textbackslash}n{\textquotedbl});}



{

}}



{

while (!appProc->isTerminated()) {}



{

bpatch.waitForStatusChange();}



{

}}



{

} else if (appBin) {}



{

if (!appBin->writeFile(newName)) {}



{

fprintf(stderr, {\textquotedbl}writeFile failed{\textbackslash}n{\textquotedbl});}



{

}}



{

}}



{

}}





\bigskip



{

int main() {}



{

// Set up information about the program to be instrumented}



{

const char* progName = {\textquotedbl}InterestingProgram{\textquotedbl};}



{

int progPID = 42;}



{

const char* progArgv[] = {{\textquotedbl}InterestingProgram{\textquotedbl}, {\textquotedbl}-h{\textquotedbl},

NULL};}



{

accessType_t mode = create;}





\bigskip



{

// Create/attach/open a binary}



{

BPatch_addressSpace* app =



{

startInstrumenting(mode, progName, progPID, progArgv);}



{

if (!app) {}



{

fprintf(stderr, {\textquotedbl}startInstrumenting failed{\textbackslash}n{\textquotedbl});}



{

exit(1);}



{

}}





\bigskip



{

// Instrument memory accesses}



{

if (!instrumentMemoryAccesses(app)) {}



{

fprintf(stderr, {\textquotedbl}instrumentMemoryAccesses failed{\textbackslash}n{\textquotedbl});}



{

exit(1);}



{

}}





\bigskip



{

// Finish instrumentation



{

const char* progName2 = {\textquotedbl}InterestingProgram-rewritten{\textquotedbl};}



{

finishInstrumenting(app, progName2);}



{

}}



